\subsection{Panel selection}

Candidate experiments were every released ENCODE eCLIP experiment in K562 and HepG2 (139 and 105
respectively). Where a protein had several peak files we took the one with the most biological
replicates, breaking ties by ENCODE's \texttt{preferred\_default} flag, and required
\texttt{output\_type=peaks}, \texttt{assembly=GRCh38}, \texttt{status=released} and
\texttt{file\_format=bed}. Every candidate was preprocessed under both composition-matching
protocols; datasets retaining at least 400 out-of-fold scored pairs were eligible.

The study panel was then defined once by the dinucleotide arm, which is the stricter control:
the 189 eligible datasets were sorted by pair count with a stable sort and every second one
retained, giving 95 datasets over 79 proteins. The set was halved rather than taken whole
because the sweep trains three model classes over five folds and three protocols, and the full
eligible set did not fit the compute budget the retained panel already consumes. Given that a
reduction was necessary, we sampled systematically by rank instead of taking the largest $N$,
because AUROC correlates with dataset size in these data, negligibly for
the composition baseline (Pearson $r$ on $\log_{10} n$ of $-0.04$ in the GC arm) but strongly
for the models, reaching $0.70$ for SpliceBERT under dinucleotide matching, so a size-truncated
panel would confound panel membership with the measured quantity; the retained panel spans the
0th to 99th percentile of the candidate size distribution. Of the 95, 94 also cleared the floor
in the GC arm and constitute the paired panel used for every three-protocol result, giving 282
protocol-dataset cells. All 95 are used only for the comparison of standalone model AUROCs
quoted in the Introduction, which needs no GC-arm counterpart. Panel membership was written once
to a manifest that every downstream stage reads. Supplementary Table S1 lists every dataset with
its ENCFF file and ENCSR experiment accession.

Sixteen proteins are assayed in both cell lines and contribute two datasets each, fifteen of
them in the paired panel. All headline intervals therefore resample proteins rather than
datasets.

\subsection{Windows}

Positive windows are 101~nt centred on the midpoint of each peak interval. Peak width and the
significance columns are unused: we applied no further $p$-value, $q$-value or fold-change
filter, because ENCODE's released \texttt{peaks} files are already thresholded and across our 95
files the minimum log2 fold-enrichment is 3.0000 and the minimum $-\log_{10}p$ is 3.0003, which
are the criteria of \citet{vannostrand2020}; 91.7\% of the 463{,}091 peaks carry ENCODE's IDR
label. The positive set is therefore every peak ENCODE called and no weaker ones. Minus-strand
windows are reverse-complemented
and transcribed to RNA, so every sequence is the strand the protein sees. Windows containing an
ambiguous base, falling outside the assembly, duplicating an existing window by start
coordinate, or carrying no region annotation are dropped and counted. Analysis is restricted to
chr1 to chr22 and chrX; chrY and all scaffolds and alternates are excluded.

Region class is assigned from GENCODE v45 \citep{frankish2023} as the first of (5$'$UTR,
3$'$UTR, CDS, non-coding exon, intron) whose interval set contains the window midpoint. Region
intervals were built from all transcripts with no
gene-type filter and merged, with strand disregarded at this step; introns were computed as the
gaps between consecutive exons.
GTF 1-based inclusive coordinates are converted to 0-based half-open.

\subsection{The three negative-set protocols}

Each protocol produces one negative per positive. All three use a single random stream seeded at
7 per protein.

\paragraph{GC-matched.} Candidates are windows of the same region class on the same chromosome,
at least 500~nt from any peak of the target protein, drawn with probability proportional to the
length of each free interval. A candidate is accepted if its GC content is within 0.05 of the
positive's. The matcher relaxes: after 25 unsuccessful draws the tolerance doubles to 0.10, and
if 40 draws fail the best candidate seen is accepted provided it is within three times the
nominal tolerance, that is 0.15. Only the target protein's own
peaks are excluded, so a GC-matched or dinucleotide-matched negative may be another protein's
binding site. The negative inherits the positive's strand label; the consequences are quantified
in the Results.

\paragraph{Dinucleotide-matched.} Positives are grouped by region and chromosome, and for each
group a candidate pool of $\max(1500,\, 8n)$ windows is sampled from the same free intervals the
GC matcher draws on: same region class, same chromosome, and at least 500~nt from any peak of the
target protein, the exclusion applied through the same code path and the same configured distance.
The two composition-matched protocols therefore share one candidate universe and differ only in
the statistic they match on. For each positive the nearest
unused candidate is assigned greedily under the $L_1$ distance over the sixteen dinucleotide
\emph{counts}, searched with a $k$-d tree over the forty nearest neighbours. The distance is over counts, not frequencies, because integer
distances are exactly representable, and each candidate carries an additional
tiebreak column that increases strictly with genomic position and is
bounded below 0.5, so it can never reorder candidates whose integer distances differ but resolves
exact ties by position. Without both devices the choice among tied candidates depends on the
floating-point behaviour of the $k$-d tree implementation and a different negative is selected
for about 8\% of pairs between CPU architectures. Assignment is greedy, not optimal;
achieved distances are reported so the quality of the approximation is visible. Used coordinates
are tracked across all groups, because one interval can be annotated as both a non-coding exon
and a 3$'$UTR and so appear in two region pools. No maximum distance is imposed.

\paragraph{Bias-aware (other RBPs' sites).} Following \citet{horlacher2023}, negatives for a
target are the positive windows of the other panel proteins assayed in the same cell line,
sampled 1:1 without replacement, excluding any donor window within 500~nt of any of the target's
own positive windows. Our version deviates from \citet{horlacher2023} in two ways: our donor
pool is the study panel (40 to 48 donor proteins per dataset, median 45) and not a full
ENCODE survey, and exclusion is by distance from the target's 101~nt windows, not from its
full peak intervals. Both the donor pool and the target's positives are read from the GC arm's
window tables, so this arm inherits the GC matcher's retained positive set and its pair count is
identical to that arm's. Sampling is performed within fold, so every pair stays inside one
cross-validation fold. These negatives are transcribed, accessible to crosslinking and
strand-correct by construction, and are not composition-matched in any respect.

\paragraph{What each protocol holds fixed.} The three protocols are not three settings of one
knob, and the difference changes how their results should be read. Both composition-matched
protocols draw candidates of the same transcript region class as the positive, so the region
label carries no information about the class. The bias-aware protocol matches fold only, which
leaves region free; how far region separates the two classes there is measured in
Section~\ref{sec:region}. Only the bias-aware
protocol's negatives are transcribed RNA, and conversely 67.6\% of GC-arm negatives lie in loci
not transcribed on the strand the window is labelled with. Section~\ref{sec:region} bounds what
the region asymmetry contributes; the strand and
transcription controls in the Results bound the other two; these are the tables the
repository names \texttt{strand\_placebo} and \texttt{transcription\_placebo}.

\subsection{Achieved match quality}

Achieved matching was quantified over every pair in each arm, 456{,}734 in the GC and
bias-aware arms and 457{,}998 in the dinucleotide arm (Table~\ref{tab:match}).

\begin{table}[htbp]
\centering
\caption{Achieved match quality between each positive and its assigned negative, all pairs.
Columns two to five describe the absolute GC difference between the pair; $L_1$ is over the
sixteen dinucleotide frequencies, on a 0 to 2 scale.}
\label{tab:match}
\begin{tabular}{lcccccc}
\toprule
arm & med.\ $|\Delta$GC$|$ & p90 & p99 & max & within 0.05 & med.\ $L_1$ \\
\midrule
GC-matched & 0.0297 & 0.0495 & 0.1089 & 0.1486 & 94.8\% & 0.500 \\
dinucleotide-matched & 0.0198 & 0.0693 & 0.1485 & 0.4159 & 85.0\% & 0.220 \\
bias-aware & 0.1387 & 0.3169 & 0.4555 & 0.6733 & 20.4\% & 0.700 \\
\bottomrule
\end{tabular}
\end{table}

The GC matcher holds its nominal tolerance for 94.8\% of pairs and its observed maximum,
0.1486, sits at the 0.15 fallback cap. ``Dinucleotide-matched'' likewise means a median $L_1$
of 0.220 rather than zero: the matcher improves composition distance 2.27-fold relative to the
GC arm but does not eliminate it. The bias-aware arm is not composition-matched at all, with a
median GC gap 4.7 times the GC arm's.

\subsection{Cross-validation}

Evaluation is grouped 5-fold cross-validation over whole chromosomes. The chromosome-to-fold map
is frozen once for all datasets in both cell lines, so the fifteen proteins measured in both
lines are evaluated on the same chromosomes and a between-line difference cannot be a partition
artefact. The map was chosen by random-restart hill climbing on peak counts only, never on labels
or model output, minimising the summed squared deviation from equal mass per fold with every
dataset weighted equally; it was required to beat round-robin assignment, to hold the median
per-dataset maximum deviation at or below 0.05, and to place no more than half of any dataset's
data in one fold. Fold $i$ is the test fold, fold $i+1 \bmod 5$ is validation and the remaining
three train, so each fold is test once, validation once and training three times, and no
additional data is spent on model selection.

Every window is scored exactly once out of fold and all AUROCs are computed on the pooled
out-of-fold score vector. Pooling, instead of averaging per-fold AUROCs, makes the
estimate invariant to fold size; the least balanced dataset places 0.42 of its data in a single
fold.

Leakage between folds was audited by exact 32-mer sharing, not by clustered sequence
identity. A 101~nt window contains 70 distinct 32-mers, and two unrelated sequences share one
with negligible probability, so a shared 32-mer indicates common origin. Over all 94 datasets,
holding fold 0 out, a median of 2.0\% of held-out windows shared any 32-mer with a training
window (maximum 24.3\%), a median of 0.08\% shared more than half of their 32-mers, and the
median dataset contained no window sharing more than 90\%. Chromosome grouping therefore leaves
a small residue of homologous sequence across folds, which is the quantity a sequence-identity
clustering step would remove.

\subsection{The composition baseline and the nested contribution}

The composition baseline is a 19-column design matrix: three mononucleotide frequencies, fifteen
dinucleotide frequencies and the Shannon entropy of the window's base composition. One level is
dropped from each frequency family because frequencies sum to one and retaining all of them
makes the design singular; dropping a level does not change the space the block spans. Every
column is standardised over the whole dataset, not within fold; the transform is
label-free and identical in both terms of every comparison, so it cannot manufacture a
contribution. GC content is not a feature: it is the spanned combination C$+$G, and
mononucleotide counts are almost, but not exactly, marginals of dinucleotide counts: over a
window of length $L$ the dinucleotide counts sum to $L-1$ and recover each base's count only up
to whether the terminal base is that base. The linear block therefore has rank 18 and not 15,
with a median 0.14\% of each mononucleotide column's variance lying outside the span of the
dinucleotide columns (maximum 0.25\% over 25 datasets), and the entropy column adds a
nineteenth, nonlinear direction. What the two protocols constrain is nonetheless asymmetric by
construction, and that asymmetry is what the design argument rests on: GC matching fixes a
single linear functional of local composition, whereas dinucleotide matching fixes all fifteen
degrees of freedom of the dinucleotide simplex, so it leaves the baseline strictly less
compositional variation to exploit. The baseline is refit on each protocol's own windows and is
never carried between protocols.

The nested contribution of a model is
\begin{equation*}
\mathrm{AUROC}(\text{composition} + \text{standardised model score}) -
\mathrm{AUROC}(\text{composition alone}).
\end{equation*}
Both terms were fitted by $L_2$-penalised logistic regression ($C=1$, lbfgs, at most 3000
iterations), once per fold on that fold's training rows, and evaluated as pooled out-of-fold
linear predictors on identical rows and folds. The model-score column is the out-of-fold score
for every row, training rows included, so the nested model is never fitted on an in-sample
score; had it been, the fitted and evaluated covariate would sit on different scales and the
mismatch would grow with how much the underlying model overfits, which differs sharply between
a 256-column penalised logistic and a fully fine-tuned transformer. The two were compared by DeLong's paired
estimator \citep{delong1988}, in the $O(n\log n)$ midrank form of \citet{sun2014}. The paired
estimator is necessary here because the two score vectors are fitted on overlapping training
data and evaluated on identical rows, so they are strongly correlated and treating them as
independent would overstate the variance of their difference: the median DeLong standard error
of the difference is 0.26 times what independence gives. Confidence intervals on a single
dataset's contribution are Wald intervals on the DeLong standard error.

\citet{demler2012} show that DeLong's test is not valid for strictly nested models, because
under the null the added predictor's contribution degenerates. Two things bound the exposure
here. The panel-level intervals, which carry every headline, do not use the DeLong standard
error at all; they are protein-clustered bootstrap intervals over per-dataset point estimates.
And where the DeLong standard error is used, for the per-dataset significance counts, it is
conservative relative to the recommended alternative: a penalised likelihood-ratio test on the
same fits calls 89 of 94 datasets significant in the GC arm against DeLong's 80. A third
feature of the design might have cut the other way and does not: negatives are matched 1:1 to
positives, so the rows are paired, while DeLong treats the two classes as independent samples.
Resampling matched pairs instead reproduces the DeLong standard error to within 1\% (median
ratio 0.99 over 30 datasets), so the matching needs no correction.

Both arms are evaluated out of fold. An in-sample composition AUROC compared against an
out-of-fold model AUROC flatters composition sufficiently to reverse the comparison on
individual datasets.

\subsection{Models}

\paragraph{4-mer logistic regression.} Raw 4-mer counts over the RNA sequence (98 per window,
256 columns) into $L_2$-penalised logistic regression ($C=1$). One model per fold, out-of-fold
score taken as the linear predictor. The $k$-mer order is 4, and the contrast is
positive at every order from 3 to 6.

\paragraph{Convolutional network.} A DeepBind-style architecture \citep{alipanahi2015}:
convolution (4 to 16 channels, width 12), ReLU, max-pool 4, convolution (16 to 32, width 8),
ReLU, global max-pool over position, dense 32 to 64, ReLU, dropout 0.5, dense 64 to 1. 7089
parameters, all trained from scratch. Global max-pooling makes the model
position-invariant, which is necessary here because exploratory analysis placed the
discriminative signal approximately 15~nt off centre, varying by protein.

\paragraph{SpliceBERT.} The pretrained nucleotide language model of \citet{chen2024},
\texttt{multimolecule/\allowbreak splicebert}, fully fine-tuned; 19.78~M parameters, all trainable. The
classification head is a mean pool over real nucleotide tokens, masking classification,
separator and padding tokens, then dense to 128, ReLU, dropout 0.3, dense to 1. Weights are
baked into the container image at build time.

Both neural models: AdamW, weight decay 0.01, batch size 32, binary cross-entropy, single
precision, at most 12 epochs with early stopping on validation AUROC at patience 4, best
checkpoint restored before test scoring. Learning rate $10^{-3}$ for the head and
$3\times10^{-5}$ for a fine-tuned encoder. Nominal epochs overstate the training
performed: the median best epoch for SpliceBERT was 3 and most runs stopped by epoch 6 or 7,
whereas the convolutional network typically ran the full 12. Initial weights were drawn before
the seed was set, so initialisation is unseeded across all 2830 committed fold-runs, that is two
models by five folds by 94 datasets in the GC and bias-aware arms and 95 in the dinucleotide
arm; the measured
cost is a per-dataset standard deviation of 0.006 to 0.010 in the nested contribution, inducing
about 0.001 on a 94-dataset mean against a reported interval half-width of about 0.008. The
convolutional and transformer arms were trained on different accelerators, which changes
floating-point summation order and is covered by the same measurement.

\subsection{Inference}

Because fifteen of the 79 proteins contribute two datasets each and the within-protein
correlation of the primary contrast is 0.92, every headline interval is a percentile interval
from a bootstrap that resamples \emph{proteins} and takes all datasets of each sampled protein,
4000 draws. Dataset-level intervals are 1.05 to 1.23 times narrower and are not reported as
headline values.

Two analyses are confirmatory and are the only ones we adjust. The first is the primary
contrast, the difference in nested contribution between protocols, which the study was built to
estimate. The second is the recommendation test of Section~\ref{sec:recommendation}, whose two
falsifying criteria were fixed before the external data were scored and whose interval is
Bonferroni-corrected over the three protocol pairs. Everything else here is exploratory. That
includes the baseline attribution, the family partition, the order-three contrast, the
within-panel reversal and the caliper-matched nulls, all of which we report as point estimates
with intervals rather than as decisions, and none of which is adjusted for the number of
analyses performed. Readers should treat an exploratory interval that just excludes zero
accordingly.

Per-dataset significance counts are also reported at a design effect of 1.35 applied to the
DeLong standard error, which is intended to cover two things the estimator does not: that it
conditions on two score vectors as fixed when both are fitted, and that it treats spatially
clustered windows as independent. We measured both components instead of assuming them
(\texttt{scripts/design\_effect.py}). Resampling whole genomic blocks inflates the standard
error of the contribution by a median factor of 1.100, and the factor is insensitive to block
length, moving only to 1.104 at 100~kb and 1.138 at 1~Mb, so the objection that a 10~kb block
is smaller than a gene does not bite. Permuting the model score within fold and refitting gives
1.048. The measured product is therefore 1.15 rather than 1.35. We retain 1.35, because it is
the more conservative of the two and because the count it produces was fixed before the
components were measured; a reader who prefers the measured figure should read 75 of 94 rather
than 72 of 94 below. This is one of two quantities in the paper that the verification harness
cannot check
from the released evidence, because the block bootstrap needs the genomic coordinates that the
published per-window score files do not carry.

\subsection{Reproducibility}

All published values are reproduced offline by a single command against committed evidence. The
harness runs 706 numeric assertions against a frozen expectations file, and the harness asserts
that no
fewer than that number ran, so a check cannot silently skip. Two assertions are
stronger than regression gates: 285 published AUROCs are rebuilt from committed per-example
scores to a maximum absolute difference of $2.2\times10^{-16}$, and the headline contrast is
rebuilt from raw sequence to $1.2\times10^{-6}$. Per-window out-of-fold scores for all three
model classes are committed, so every cell of the model-class comparison is recomputable from the
repository alone, and is not asserted against a hand-entered table.

Three devices guard against platform-dependent results: dinucleotide distances are computed on
integer counts, candidate ties are broken by genomic position, and the panel sort is stable.
Each was added after an observed discrepancy between CPU architectures or between runs.

Analyses used NumPy \citep{numpy}, SciPy \citep{scipy}, scikit-learn \citep{sklearn}, PyTorch
\citep{pytorch} and Matplotlib \citep{matplotlib}.
